%% sections/03-background.tex
%% Background & Theory
%% IEC/IEEE 82079-1:2019 — Conceptual Information (§5)
%% ─────────────────────────────────────────────────────

\section{Background and Theory}
\label{sec:background}

%%
%% WRITING GUIDE — BACKGROUND:
%%
%%   Purpose: Give the reader enough conceptual context to understand WHY
%%   each step in Section 6 is performed.  Do NOT repeat the procedure here.
%%
%%   Subsections you typically need:
%%     3.1  System overview (what the machine does in this mode)
%%     3.2  Key equations (force model, sensor transfer function)
%%     3.3  Control loop architecture
%%     3.4  Relevant ACS axis/parameter names
%%
%%   Keep equations self-contained: define every symbol.
%%   Use \cite{key} for any external reference (journal paper, datasheet).
%%

\subsection{System Overview}
\label{sec:background:system}

The PBA GS-Series gantry operates in either position control or force
control mode.  In \emph{force control mode}, the ACS SPiiPlus controller
replaces the position error with a force error signal derived from an
analogue load cell connected to the controller's analogue input port (AIN).
The control architecture is illustrated in \cref{fig:control-arch}.

%%
%% FIGURE — Replace the placeholder below with your actual architecture diagram.
%% Place the image file in the /figures/ folder.
%%
\begin{figure}[H]
  \centering
  %% \includegraphics[width=0.85\textwidth]{figures/force-control-arch.png}
  \fbox{\parbox{0.85\textwidth}{\centering\vspace{20mm}
    \textcolor{PBAGray}{[FIGURE PLACEHOLDER: Force control architecture diagram]}\vspace{20mm}}}
  \caption{Force control closed-loop architecture for PBA GS-Series gantry.
           The load cell signal is conditioned and fed to AIN~1 of the
           ACS SPiiPlus controller.}
  \label{fig:control-arch}
\end{figure}


\subsection{Force–Displacement Model}
\label{sec:background:model}

The relationship between Z-axis motor current and applied contact force is
approximated as linear within the operating range:

\begin{equation}
  F_{\text{applied}} = k_f \cdot I_{\text{cmd}} - F_{\text{preload}}
  \label{eq:force-model}
\end{equation}

\noindent where:
\begin{itemize}[noitemsep, topsep=4pt]
  \item $F_{\text{applied}}$ is the force applied to the workpiece [\si{\newton}]
  \item $k_f$ is the force constant of the linear motor [\si{\newton\per\ampere}]
  \item $I_{\text{cmd}}$ is the commanded current [\si{\ampere}]
  \item $F_{\text{preload}}$ is the gravitational preload of the Z-axis assembly [\si{\newton}]
\end{itemize}

The load cell provides a feedback signal $V_{\text{cell}}$ [\si{\volt}] that is
converted to force units using the calibrated sensitivity $S$ [\si{\newton\per\volt}]:

\begin{equation}
  F_{\text{measured}} = S \cdot (V_{\text{cell}} - V_{\text{offset}})
  \label{eq:sensor-model}
\end{equation}

For typical Kistler~9317C sensors used on PBA systems, $S \approx
\SI{100}{\newton\per\volt}$; confirm with your specific sensor datasheet
\cite{kistler_sensor_2022}.


\subsection{PID Force Loop}
\label{sec:background:pid}

The ACS SPiiPlus force loop is implemented as a standard PID controller
operating at the servo update rate:

\begin{equation}
  u(t) = K_p \, e(t) + K_i \int_0^t e(\tau)\,\mathrm{d}\tau
         + K_d \,\frac{\mathrm{d}e(t)}{\mathrm{d}t}
  \label{eq:pid}
\end{equation}

\noindent where $e(t) = F_{\text{ref}} - F_{\text{measured}}$ is the force error.

The controller gains \acsparam{KP}, \acsparam{KI}, and \acsparam{KD} are
set in the ACS parameter table (see \cref{sec:parameters}).
Initial values from the PBA commissioning baseline are given in
\cref{tab:pid-defaults}.

\begin{table}[H]
  \centering
  \caption{Baseline PID gains for force control on PBA GS-Series.
           Adjust \acsparam{KP} first, then \acsparam{KI}.}
  \label{tab:pid-defaults}
  \begin{tabular}{@{} l l l l @{}}
    \toprule
    \textbf{Parameter} & \textbf{Description} & \textbf{Baseline value} & \textbf{Unit} \\
    \midrule
    \acsparam{KP} & Proportional gain & 2.5 & \si{\ampere\per\newton} \\
    \rowcolor{PBALightGray}
    \acsparam{KI} & Integral gain & 0.08 & \si{\ampere\per\newton\per\second} \\
    \acsparam{KD} & Derivative gain & 0.001 & \si{\ampere\second\per\newton} \\
    \rowcolor{PBALightGray}
    \acsparam{KP\_FFW} & Velocity feedforward gain & 0.0 & --- \\
    \bottomrule
  \end{tabular}
\end{table}

\begin{notebox}
  The baseline gains above are starting points derived from PBA internal
  testing.  Actual optimal gains depend on workpiece stiffness, Z-axis mass,
  and contact geometry.  Always tune on-site using the verification procedure
  in \cref{sec:verification}.
\end{notebox}


\subsection{Key References}
\label{sec:background:refs}

Background reading relevant to this setup:
force sensor selection and conditioning \cite{kistler_sensor_2022},
ACS SPiiPlus motion controller architecture \cite{acs_spiiplus_ref},
and general force control for precision bonding \cite{compliant_motion_2018}.
